\documentclass[10pt,a4paper]{amsart}
\usepackage{amsmath,amssymb,amsthm,graphicx}
% Retained preamble for the Stacks Project chapter snapshot in this directory.
% The chapter's own preamble.tex is a copy of the shared Stacks Project preamble; the block
% below is sliced MECHANICALLY, line for line, from it. Nothing was added, removed or
% redefined; the wrapper class, the xr-hyper/hyperref wiring and the externaldocument list
% are omitted because the assembler supplies the document class.
%
% lines taken from preamble.tex: 3, 5, 9, 11, 12, 156, 157, 158, ... (51 lines total)
%
% package coverage restored from the chapter's own preamble (2026-10-13): xy, xy-2cell

\usepackage[all]{xy}
\xyoption{2cell}
\UseAllTwocells

\newenvironment{reference}{\comment}{\endcomment}
\newenvironment{slogan}{\comment}{\endcomment}
\newenvironment{history}{\comment}{\endcomment}
\theoremstyle{plain}
\newtheorem{theorem}[subsection]{Theorem}
\newtheorem{proposition}[subsection]{Proposition}
\newtheorem{lemma}[subsection]{Lemma}
\theoremstyle{definition}
\newtheorem{definition}[subsection]{Definition}
\newtheorem{example}[subsection]{Example}
\newtheorem{exercise}[subsection]{Exercise}
\newtheorem{situation}[subsection]{Situation}
\theoremstyle{remark}
\newtheorem{remark}[subsection]{Remark}
\newtheorem{remarks}[subsection]{Remarks}
\numberwithin{equation}{subsection}
\def\lim{\mathop{\mathrm{lim}}\nolimits}
\def\colim{\mathop{\mathrm{colim}}\nolimits}
\def\Spec{\mathop{\mathrm{Spec}}}
\def\Hom{\mathop{\mathrm{Hom}}\nolimits}
\def\Ext{\mathop{\mathrm{Ext}}\nolimits}
\def\SheafHom{\mathop{\mathcal{H}\!\mathit{om}}\nolimits}
\def\SheafExt{\mathop{\mathcal{E}\!\mathit{xt}}\nolimits}
\def\Sch{\mathit{Sch}}
\def\Mor{\mathop{\mathrm{Mor}}\nolimits}
\def\Ob{\mathop{\mathrm{Ob}}\nolimits}
\def\Sh{\mathop{\mathit{Sh}}\nolimits}
\def\NL{\mathop{N\!L}\nolimits}
\def\CH{\mathop{\mathrm{CH}}\nolimits}
\def\proetale{{pro\text{-}\acute{e}tale}}
\def\etale{{\acute{e}tale}}
\def\QCoh{\mathit{QCoh}}
\def\Ker{\mathop{\mathrm{Ker}}}
\def\Im{\mathop{\mathrm{Im}}}
\def\Coker{\mathop{\mathrm{Coker}}}
\def\Coim{\mathop{\mathrm{Coim}}}
\DeclareMathSymbol{\boxtimes}{\mathbin}{AMSa}{"02}
\def\QCohstack{\mathcal{QC}\!\mathit{oh}}
\def\Cohstack{\mathcal{C}\!\mathit{oh}}
\def\Spacesstack{\mathcal{S}\!\mathit{paces}}
\def\Quotfunctor{\mathrm{Quot}}
\def\Hilbfunctor{\mathrm{Hilb}}
\def\Curvesstack{\mathcal{C}\!\mathit{urves}}
\def\Polarizedstack{\mathcal{P}\!\mathit{olarized}}
\def\Complexesstack{\mathcal{C}\!\mathit{omplexes}}
\def\Pic{\mathop{\mathrm{Pic}}\nolimits}
\def\Picardstack{\mathcal{P}\!\mathit{ic}}
\def\Picardfunctor{\mathrm{Pic}}
\def\Deformationcategory{\mathcal{D}\!\mathit{ef}}
\usepackage{cleveref}
\crefname{originalreference}{original reference}{original references}
\Crefname{originalreference}{Original reference}{Original references}
\makeatletter
\newlabel{cohomology-remark-base-change}{{08HY}{source}{}{sourcepointer-1}{}}
\newlabel{cohomology-remark-base-change@cref}{{[originalreference][0][]08HY}{[1][0][]source}{}{}{}}
\makeatother
\crefname{originalreference}{original reference}{original references}
\Crefname{originalreference}{Original reference}{Original references}
\makeatletter
\newlabel{properties-lemma-extend-trivial}{{01PE}{source}{}{sourcepointer-2}{}}
\newlabel{properties-lemma-extend-trivial@cref}{{[originalreference][0][]01PE}{[1][0][]source}{}{}{}}
\makeatother
\crefname{originalreference}{original reference}{original references}
\Crefname{originalreference}{Original reference}{Original references}
\makeatletter
\newlabel{coherent-lemma-relative-affine-vanishing}{{01XC}{source}{}{sourcepointer-3}{}}
\newlabel{coherent-lemma-relative-affine-vanishing@cref}{{[originalreference][0][]01XC}{[1][0][]source}{}{}{}}
\makeatother
\crefname{originalreference}{original reference}{original references}
\Crefname{originalreference}{Original reference}{Original references}
\makeatletter
\newlabel{schemes-lemma-widetilde-pullback}{{01I9}{source}{}{sourcepointer-4}{}}
\newlabel{schemes-lemma-widetilde-pullback@cref}{{[originalreference][0][]01I9}{[1][0][]source}{}{}{}}
\makeatother
\crefname{originalreference}{original reference}{original references}
\Crefname{originalreference}{Original reference}{Original references}
\makeatletter
\newlabel{coherent-lemma-quasi-coherence-higher-direct-images-application}{{01XK}{source}{}{sourcepointer-5}{}}
\newlabel{coherent-lemma-quasi-coherence-higher-direct-images-application@cref}{{[originalreference][0][]01XK}{[1][0][]source}{}{}{}}
\makeatother
\crefname{originalreference}{original reference}{original references}
\Crefname{originalreference}{Original reference}{Original references}
\makeatletter
\newlabel{cohomology-lemma-quasi-separated-cohomology-colimit}{{01FF}{source}{}{sourcepointer-6}{}}
\newlabel{cohomology-lemma-quasi-separated-cohomology-colimit@cref}{{[originalreference][0][]01FF}{[1][0][]source}{}{}{}}
\makeatother
\crefname{originalreference}{original reference}{original references}
\Crefname{originalreference}{Original reference}{Original references}
\makeatletter
\newlabel{coherent-lemma-quasi-coherence-higher-direct-images}{{01XJ}{source}{}{sourcepointer-7}{}}
\newlabel{coherent-lemma-quasi-coherence-higher-direct-images@cref}{{[originalreference][0][]01XJ}{[1][0][]source}{}{}{}}
\makeatother
\crefname{originalreference}{original reference}{original references}
\Crefname{originalreference}{Original reference}{Original references}
\makeatletter
\newlabel{coherent-lemma-relative-affine-cohomology}{{089W}{source}{}{sourcepointer-8}{}}
\newlabel{coherent-lemma-relative-affine-cohomology@cref}{{[originalreference][0][]089W}{[1][0][]source}{}{}{}}
\makeatother
\title{Cohomological dimension d controls base change of the top direct image}
\author{The Stacks Project Authors}
\date{stacks-project:limits.tex:local-snapshot-20261012; GNU FDL 1.2 or later}
\begin{document}
\maketitle
\noindent\textit{Editorial note.} Verbatim author omission in the counted span: 'Any right exact R'-linear functor F which commutes with direct sums is given by tensoring with an R'-module (omitted; left as exercise for the reader).' This is the named Eilenberg-Watts theorem, an external established result; the load-bearing citation is to a classic, and (a)/(c) are fully worked in span. Flagged needs\_human\_spotcheck for exactly this sentence.\par\medskip
\noindent\textit{Allowed standard prerequisites (editorial).}\begin{itemize}
\item Eilenberg-Watts tensor representation. classic theorem; the source marks the step '(omitted; left as exercise for the reader)'
\end{itemize}\medskip
\ifdefined\hypertarget\hypertarget{sourcepointer-1}{}\fi
\noindent\textit{Source reference.} Original reference 08HY: sources/stacks-chow/tags.txt, source lines 10796--10796. Cross-chapter standard prerequisite cited by the author inside the counted span; it is used as a pointer and its own proof is not part of this item.\par\medskip
\ifdefined\hypertarget\hypertarget{sourcepointer-2}{}\fi
\noindent\textit{Source reference.} Original reference 01PE: sources/stacks-chow/tags.txt, source lines 2227--2227. Cross-chapter standard prerequisite cited by the author inside the counted span; it is used as a pointer and its own proof is not part of this item.\par\medskip
\ifdefined\hypertarget\hypertarget{sourcepointer-3}{}\fi
\noindent\textit{Source reference.} Original reference 01XC: sources/stacks-chow/tags.txt, source lines 2513--2513. Cross-chapter standard prerequisite cited by the author inside the counted span; it is used as a pointer and its own proof is not part of this item.\par\medskip
\ifdefined\hypertarget\hypertarget{sourcepointer-4}{}\fi
\noindent\textit{Source reference.} Original reference 01I9: sources/stacks-chow/tags.txt, source lines 1970--1970. Cross-chapter standard prerequisite cited by the author inside the counted span; it is used as a pointer and its own proof is not part of this item.\par\medskip
\ifdefined\hypertarget\hypertarget{sourcepointer-5}{}\fi
\noindent\textit{Source reference.} Original reference 01XK: sources/stacks-chow/tags.txt, source lines 2521--2521. Cross-chapter standard prerequisite cited by the author inside the counted span; it is used as a pointer and its own proof is not part of this item.\par\medskip
\ifdefined\hypertarget\hypertarget{sourcepointer-6}{}\fi
\noindent\textit{Source reference.} Original reference 01FF: sources/stacks-chow/tags.txt, source lines 1868--1868. Cross-chapter standard prerequisite cited by the author inside the counted span; it is used as a pointer and its own proof is not part of this item.\par\medskip
\ifdefined\hypertarget\hypertarget{sourcepointer-7}{}\fi
\noindent\textit{Source reference.} Original reference 01XJ: sources/stacks-chow/tags.txt, source lines 2520--2520. Cross-chapter standard prerequisite cited by the author inside the counted span; it is used as a pointer and its own proof is not part of this item.\par\medskip
\ifdefined\hypertarget\hypertarget{sourcepointer-8}{}\fi
\noindent\textit{Source reference.} Original reference 089W: sources/stacks-chow/tags.txt, source lines 10511--10511. Cross-chapter standard prerequisite cited by the author inside the counted span; it is used as a pointer and its own proof is not part of this item.\par\medskip
\begin{lemma}
\label{lemma-top-cohomology-functor}
Let $f : X \to Y$ be a morphism of schemes. Let $d \geq 0$. Assume
\begin{enumerate}
\item $X$ and $Y$ are quasi-compact and quasi-separated, and
\item $R^if_*\mathcal{F} = 0$ for $i > d$ and
every quasi-coherent $\mathcal{O}_X$-module $\mathcal{F}$.
\end{enumerate}
Then we have
\begin{enumerate}
\item[(a)] for any base change diagram
$$
\xymatrix{
X' \ar[d]_{f'} \ar[r]_{g'} & X \ar[d]^f \\
Y' \ar[r]^g & Y
}
$$
we have $R^if'_*\mathcal{F}' = 0$ for $i > d$ and any quasi-coherent
$\mathcal{O}_{X'}$-module $\mathcal{F}'$,
\item[(b)]
$R^df'_*(\mathcal{F}' \otimes_{\mathcal{O}_{X'}} (f')^*\mathcal{G}') =
R^df'_*\mathcal{F}' \otimes_{\mathcal{O}_{Y'}} \mathcal{G}'$
for any quasi-coherent $\mathcal{O}_{Y'}$-module $\mathcal{G}'$,
\item[(c)] formation of $R^df'_*\mathcal{F}'$ commutes with arbitrary
further base change (see proof for explanation).
\end{enumerate}
\end{lemma}

\begin{proof}
Before giving the proofs, we explain the meaning of (c). Suppose
we have an additional cartesian square
$$
\xymatrix{
X'' \ar[d]_{f''} \ar[r]_{h'} &
X' \ar[d]_{f'} \ar[r]_{g'} &
X \ar[d]^f \\
Y'' \ar[r]^h &
Y' \ar[r]^g &
Y
}
$$
tacked onto our given diagram. If (a) holds, then there is a canonical
map $\gamma : h^*R^df'_*\mathcal{F}' \to R^df''_*(h')^*\mathcal{F}'$.
Namely, $\gamma$ is the map on degree $d$ cohomology sheaves
induced by the composition
$$
Lh^*Rf'_*\mathcal{F}' \longrightarrow
Rf''_*L(h')^*\mathcal{F}' \longrightarrow Rf''_*(h')^*\mathcal{F}'
$$
Here the first arrow is the base change map
(Cohomology, Remark \ref{cohomology-remark-base-change}) and
the second arrow complex from the canonical map
$L(g')^*\mathcal{F} \to (g')^*\mathcal{F}$.
Similarly, since $Rf'_*\mathcal{F}$ has no nonzero cohomology sheaves in
degrees $> d$ by (a) we have
$H^d(Lh^*Rf_*\mathcal{F}') = h^*R^df_*\mathcal{F}$.
The content of (c) is that $\gamma$ is an isomorphism.

\medskip\noindent
Having said this, we can check (a), (b), and (c) locally on $Y'$ and $Y''$.
Suppose that $V \subset Y$ is a quasi-compact open subscheme. Then we
claim (1) and (2) hold for $f|_{f^{-1}(V)} : f^{-1}(V) \to V$.
Namely, (1) is immediate and (2) follows because any quasi-coherent
module on $f^{-1}(V)$ is the restriction of a quasi-coherent module
on $X$ (Properties, Lemma \ref{properties-lemma-extend-trivial}) and formation
of higher direct images commutes with restriction to opens.
Thus we may also work locally on $Y$. In other words, we may
assume $Y''$, $Y'$, and $Y$ are affine schemes.

\medskip\noindent
Proof of (a) when $Y'$ and $Y$ are affine. In this case the morphisms
$g$ and $g'$ are affine. Thus $g_* = Rg_*$ and $g'_* = Rg'_*$
(Cohomology of Schemes, Lemma \ref{coherent-lemma-relative-affine-vanishing})
and $g_*$ is identified with the restriction functor on modules
(Schemes, Lemma \ref{schemes-lemma-widetilde-pullback}).
Then
$$
g_*(R^if'_*\mathcal{F}') = H^i(Rg_*Rf'_*\mathcal{F}') =
H^i(Rf_*Rg'_*\mathcal{F}') =
H^i(Rf_*g'_*\mathcal{F}') =
Rf^i_*g'_*\mathcal{F}'
$$
which is zero by assumption (2). Hence (a) by our description of $g_*$.

\medskip\noindent
Proof of (b) when $Y'$ is affine, say $Y' = \Spec(R')$.
By part (a) we have $H^{d + 1}(X', \mathcal{F}') = 0$
for any quasi-coherent $\mathcal{O}_{X'}$-module $\mathcal{F}'$, see
Cohomology of Schemes, Lemma
\ref{coherent-lemma-quasi-coherence-higher-direct-images-application}.
Consider the functor $F$ on $R'$-modules defined by the rule
$$
F(M) = H^d(X', \mathcal{F}' \otimes_{\mathcal{O}_{X'}} (f')^*\widetilde{M})
$$
By Cohomology, Lemma \ref{cohomology-lemma-quasi-separated-cohomology-colimit}
this functor commutes with direct sums (this is where we use that
$X$ and hence $X'$ is quasi-compact and quasi-separated).
On the other hand, if $M_1 \to M_2 \to M_3 \to 0$ is an exact sequence,
then
$$
\mathcal{F}' \otimes_{\mathcal{O}_{X'}} (f')^*\widetilde{M}_1 \to
\mathcal{F}' \otimes_{\mathcal{O}_{X'}} (f')^*\widetilde{M}_2 \to
\mathcal{F}' \otimes_{\mathcal{O}_{X'}} (f')^*\widetilde{M}_3 \to 0
$$
is an exact sequence of quasi-coherent modules on $X'$
and by the vanishing of higher cohomology given above we get an exact sequence
$$
F(M_1) \to F(M_2) \to F(M_3) \to 0
$$
In other words, $F$ is right exact. Any right exact $R'$-linear functor
$F : \text{Mod}_{R'} \to \text{Mod}_{R'}$
which commutes with direct sums is given by tensoring with an $R'$-module
(omitted; left as exercise for the reader).
Thus we obtain $F(M) = H^d(X', \mathcal{F}') \otimes_{R'} M$.
Since $R^d(f')_*\mathcal{F}'$ and
$R^d(f')_*(\mathcal{F}' \otimes_{\mathcal{O}_{X'}} (f')^*\widetilde{M})$
are quasi-coherent (Cohomology of Schemes, Lemma
\ref{coherent-lemma-quasi-coherence-higher-direct-images}),
the fact that $F(M) = H^d(X', \mathcal{F}') \otimes_{R'} M$
translates into the statement given in (b).

\medskip\noindent
Proof of (c) when $Y'' \to Y' \to Y$ are morphisms of affine schemes.
Say $Y'' = \Spec(R'')$ and $Y' = \Spec(R')$.
Then we see that $R^df''_*(h')^*\mathcal{F}'$
is the quasi-coherent module on $Y'$ associated to the $R''$-module
$H^d(X'', (h')^*\mathcal{F}')$. Now $h' : X'' \to X'$ is affine
hence $H^d(X'', (h')^*\mathcal{F}') = H^d(X, h'_*(h')^*\mathcal{F}')$
by the already used
Cohomology of Schemes, Lemma \ref{coherent-lemma-relative-affine-cohomology}.
We have
$$
h'_*(h')^*\mathcal{F}' =
\mathcal{F}' \otimes_{\mathcal{O}_{X'}} (f')^*\widetilde{R''}
$$
as the reader sees by checking on an affine open covering. Thus
$H^d(X'', (h')^*\mathcal{F}') = H^d(X', \mathcal{F}') \otimes_{R'} R''$
by part (b) applied to $f'$ and the proof is complete.
\end{proof}

\end{document}
